\documentclass[11pt]{article}

\usepackage[T1]{fontenc}
\usepackage[a4paper,margin=2.7cm]{geometry}
\usepackage{amsmath,amssymb,amsthm,mathtools}
\usepackage{xcolor}
\usepackage{hyperref}
\usepackage{microtype}

\hypersetup{colorlinks=true,linkcolor=blue!55!black,citecolor=blue!55!black,
  urlcolor=blue!55!black}

\newtheorem{theorem}{Theorem}[section]
\newtheorem{lemma}[theorem]{Lemma}
\newtheorem{proposition}[theorem]{Proposition}
\newtheorem{corollary}[theorem]{Corollary}
\theoremstyle{definition}
\newtheorem{definition}[theorem]{Definition}
\theoremstyle{remark}
\newtheorem*{remark}{Remark}

\newcommand{\C}{\mathbb{C}}
\newcommand{\Q}{\mathbb{Q}}
\newcommand{\Qbar}{\overline{\Q}}
\newcommand{\hull}{\operatorname{hull}}

\title{Conjugation is not independent data:\\
a formalized obstruction for Diaz's modulus question}
\author{Carlo Perassi}
\date{}

\begin{document}
\maketitle

\begin{abstract}
Let $K\subseteq\C$ be a conjugation-stable subfield and let $u\in\C$ satisfy
$u\bar u\in K$. We record an elementary observation and its consequences,
all machine-checked in Lean~4 with Mathlib: since $\bar u=(u\bar u)/u$, complex
conjugation on the subfield generated by $K$ and $u$ is determined by the ring
structure rather than being extra data. Three consequences follow. Any ring
homomorphism fixing $K$ and sending $u$ to a second point $t$ on the same circle
automatically intertwines conjugation on that subfield; such a homomorphism
exists on the subfield whenever $u$ and $t$ are transcendental over $K$; and no
vanishing-coefficient statement over $K$ built from the associated $2\times2$
matrix distinguishes $u$ from $t$. These are negative results about a class of
proof strategies. No new transcendence result is proved, and Diaz's question
remains open.
\end{abstract}

\section{Setting}

Write $\Qbar$ for the field of algebraic numbers in $\C$ and
\[
  \mathcal{L}=\{u\in\C:\ e^{u}\in\Qbar^{\times}\}
\]
for the logarithms of algebraic numbers. Diaz~\cite{Diaz2004} studied the use of
complex conjugation in transcendence arguments for the exponential function,
asking in particular how the non-holomorphic maps $z\mapsto\bar z$ and
$z\mapsto|z|$ could enter such a proof at all.

\begin{remark}[to be completed]
The precise statement and location of the modulus question --- that no non-zero
element of $\mathcal{L}$ has algebraic modulus --- must be quoted from
\cite{Diaz2004} with its section or page number before this note is circulated.
Nothing below depends on the attribution: the results are proved for an
arbitrary conjugation-stable $K$ and are stated that way.
\end{remark}

Throughout, $K\subseteq\C$ is a subfield with $\bar a\in K$ for every $a\in K$.
The case of interest is $K=\Qbar$.

\begin{definition}
For $K$ a subfield of $\C$ and $u\in\C$, let $\hull_K(u)$ denote the subfield of
$\C$ generated by $K$ together with $u$.
\end{definition}

Call $u$ a \emph{candidate over $K$} if $u\bar u\in K$; the intended case is
$u\in\mathcal{L}$ with $|u|$ algebraic, so that $u\bar u=|u|^2\in\Qbar$.

\section{The mechanism}

The whole note rests on one identity.

\begin{lemma}\label{lem:rho}
If $u\neq0$ and $\rho=u\bar u$, then $\bar u=\rho/u$.
\end{lemma}

\begin{proof}
Immediate from $\rho = u\bar u$ and $u \neq 0$.
\end{proof}

Trivial as it is, this is the point. For a candidate, $\rho$ lies in $K$, so
$\bar u$ is a rational function of $u$ with coefficients in $K$. Conjugation on
$\hull_K(u)$ is therefore not additional structure that a homomorphism might
fail to respect; it is already fixed by the field operations.

\section{Rigidity}

\begin{lemma}\label{lem:eqon}
Let $f,g\colon\C\to\C$ be ring homomorphisms agreeing on $K$ and at $u$. Then
$f=g$ on $\hull_K(u)$.
\end{lemma}

\begin{theorem}\label{thm:conjcomm}
Let $u,t\neq0$ with $u\bar u\in K$ and $t\bar t=u\bar u$. If
$\Phi\colon\C\to\C$ is a ring homomorphism fixing $K$ pointwise with
$\Phi(u)=t$, then
\[
  \Phi(\bar z)=\overline{\Phi(z)}\qquad\text{for every }z\in\hull_K(u).
\]
\end{theorem}

\begin{proof}
Both $\Phi\circ{}$conj and conj${}\circ\Phi$ are ring homomorphisms agreeing on
$K$, because $K$ is conjugation-stable and $\Phi$ fixes it. At $u$ they agree
too: by Lemma~\ref{lem:rho}, $\Phi(\bar u)=\Phi(\rho/u)=\rho/t$ and
$\bar t=\rho/t$. Now apply Lemma~\ref{lem:eqon}.
\end{proof}

The intertwining is not arranged. It is forced.

\section{Existence}

Theorem~\ref{thm:conjcomm} is conditional on a homomorphism existing. On the
hull, one does.

\begin{theorem}\label{thm:exists}
If $u$ and $t$ are both transcendental over $K$, there is a $K$-algebra
homomorphism $\Phi\colon K(u)\to\C$ with $\Phi(u)=t$.
\end{theorem}

\begin{proof}
For $f$ transcendental over $K$, the field $K(f)$ is $K$-isomorphic to the field
of rational functions $K(X)$, by $X\mapsto f$. Applying this to $u$ and to $t$
gives $K(u)\cong K(X)\cong K(t)$ over $K$; compose with the inclusion
$K(t)\hookrightarrow\C$.
\end{proof}

\begin{remark}[scope of Theorem~\ref{thm:exists}]
The homomorphism is produced on $K(u)$, not on all of $\C$. Extending a
homomorphism from $K(u)$ to an endomorphism of $\C$ is possible, by extending a
transcendence basis and using that $\C$ is algebraically closed, but it is not
proved here and is not needed: every use below evaluates $\Phi$ only on
$\hull_K(u)$.
\end{remark}

\section{The obstruction}

Attach to a candidate the matrix
\[
  H(u,r)=\begin{pmatrix} u & r\\ r & \bar u\end{pmatrix},
\]
where $r\in K$ satisfies $u\bar u=r^{2}$. Then $\det H=u\bar u-r^{2}=0$, so $H$
has rank one, and its coefficients factor.

\begin{theorem}\label{thm:novanish}
Let $r\in K$, let $u\notin K$ with $u\bar u=r^{2}$, and let
$w,v\in K^{2}$ be non-zero. Then
\[
  \sum_{i,j} w_i\,H(u,r)_{ij}\,v_j \neq 0 .
\]
\end{theorem}

\begin{corollary}\label{cor:indist}
Let $r\in K$ and let $u,t\notin K$ both satisfy $u\bar u=t\bar t=r^{2}$. Then
for all non-zero $w,v\in K^{2}$,
\[
  \sum_{i,j} w_i\,H(u,r)_{ij}\,v_j = 0
  \iff
  \sum_{i,j} w_i\,H(t,r)_{ij}\,v_j = 0 .
\]
\end{corollary}

\begin{proof}
By Theorem~\ref{thm:novanish} neither side ever holds.
\end{proof}

Corollary~\ref{cor:indist} is the negative statement. A proof strategy that
accumulates vanishing-coefficient constraints over $K$ on $u$ and $\bar u$,
hoping they collide, cannot separate a candidate from an ordinary point placed
on the same circle. The reason is Lemma~\ref{lem:rho}: the constraints are not
independent, because $\bar u$ was never independent of $u$.

\section{What is and is not claimed}

Claimed: the statements above, for an arbitrary conjugation-stable subfield $K$,
each machine-checked.

Not claimed. No transcendence result, and no progress on whether a non-zero
element of $\mathcal{L}$ can have algebraic modulus; that question is untouched
and remains open. Nothing here shows the conjugation method fails in general ---
only that this particular family of statements, vanishing of coefficients over
$K$ of the rank-one matrix $H$, does not separate candidates. Theorem
\ref{thm:exists} is proved on $K(u)$ only. The formalization does not assert
that any candidate exists; establishing that a non-zero $u\in\mathcal{L}$ with
$|u|$ algebraic is transcendental requires Hermite--Lindemann, which the
compared surface of this note deliberately avoids.

\section{Related work and search}

\begin{remark}[to be completed by the author]
This section must record an explicit literature search: which databases and
review services were consulted, with what queries and on what date; which
adjacent results were examined; and what was concluded about whether the
observations above are new. If the search does not settle novelty, this section
must say that novelty is unknown rather than assert priority. Author's existing
search notes and references are to be inserted here.
\end{remark}

\section{Formalization}

The development is Lean~4 with Mathlib; sources at
\url{https://github.com/carlok/diaz-modulus-lean}. The advertised statements are
in \texttt{Challenge.lean}, with proofs in \texttt{Solution.lean}.

\begin{center}
\begin{tabular}{lll}
\hline
Statement & Lean declaration & Module\\
\hline
Lemma~\ref{lem:rho} & \texttt{conj\_eq\_rho\_div} & \texttt{Diaz.Closure}\\
Lemma~\ref{lem:eqon} & \texttt{eqOn\_hull} & \texttt{Diaz.Closure}\\
Theorem~\ref{thm:conjcomm} & \texttt{conj\_comm} & \texttt{Diaz.Transfer}\\
Theorem~\ref{thm:exists} & \texttt{exists\_algHom\_of\_transcendental} & \texttt{Diaz.Palomar}\\
Theorem~\ref{thm:novanish} & \texttt{no\_vanishing\_coeff\_matrix} & \texttt{Diaz.Rigidity}\\
Corollary~\ref{cor:indist} & \texttt{coeff\_indistinguishable} & \texttt{Diaz.Palomar}\\
$\hull_K(u)$ & \texttt{hull} & \texttt{Diaz.Closure}\\
$H(u,r)$ & \texttt{Hmat} & \texttt{Diaz.Rigidity}\\
\hline
\end{tabular}
\end{center}

Running \verb|#print axioms| on each of these lists only \texttt{propext},
\texttt{Classical.choice} and \texttt{Quot.sound}. The wider repository also
contains two imported transcendence statements declared as axioms,
\texttt{hermite\_lindemann} and \texttt{exists\_ringHom\_of\_transcendental},
used for results outside the list above. No statement in the table depends on
either.

\begin{thebibliography}{9}

\bibitem{Diaz2004}
G.~Diaz,
\emph{Utilisation de la conjugaison complexe dans l'\'etude de la transcendance
de valeurs de la fonction exponentielle usuelle},
J.~Th\'eor. Nombres Bordeaux \textbf{16} (2004), no.~3, 535--553,
\href{https://doi.org/10.5802/jtnb.459}{doi:10.5802/jtnb.459}.

\bibitem{mathlib}
The Mathlib Community,
\emph{Mathlib: the mathematical library of Lean~4},
\href{https://github.com/leanprover-community/mathlib4}{project repository}.

\end{thebibliography}

\end{document}
